\section{Related Work}

The section provides the review of the literature on the subject of the suggested framework, analysing the previous studies in the field of demand forecasting, feature prioritisation, time series analysis, and sentiment analysis. The review reveals the gaps in existing solutions, which this study fills.

\subsection{Software Feature Demand Forecasting}
The conventional research in demand forecasting started in the setting of physical products---consumer goods, electronics and retail inventory. Box and Jenkins models of ARIMA are still in use today as they are robust and can be interpreted. Nevertheless, software capabilities have different peculiarities that distinguish the product of software, in comparison with tangible goods. The demand for features is not subjected to any predictable seasonal pattern and can fluctuate fast, in response to the announcements of viral content, rivalry, or changes in user behaviour. It has been established that the trends in internet search are significantly correlated with real-life patterns of interest. Search data has been used to forecast phenomena as diverse as financial market trends, and epidemiological trends. This theoretical reasoning has been supported by the fact that search behaviour as an indicator of true user interest, is the premise upon which publicly available search data are used as a proxy of demand. It is an understanding that motivated the use of Google Trends data as one of the main demand signals in this model.

\subsection{Prioritisation of Features on SaaS Products}
The product management practitioners use many prioritisation models such as RICE scoring, Kano models and value-versus-effort matrices. Although these frameworks offer formal systems of decision making, they need considerable subjectiveness. Practitioners need to approximate the effect, gauge user worth, and prioritise the features according to in-house deliberations. These methods are sufficient in the case of simple products, but they scale poorly and are prone to numerous types of cognitive biases. There are other organisations that have implemented analytics-based prioritisation approaches that use internal usage data to make roadmap decisions. Nonetheless, this internal data is owned and not available to any external researchers, so it is not possible to study, replicate, and generalise results across products. This study solves this weakness by showing that publicly accessible data has enough signal strength to be used to make decisions on the prioritisation of features.

\subsection{Time Series and Machine Learning}
Methods of forecasting are across a continuum between traditional statistical processes and modern machine learning techniques. Traditional methods such as Exponential Smoothing and ARIMA models are easily interpretable, have a solid theoretical basis, and do not require extensive data. Neural networks, gradient boosting, random forests, the various machine learning methods, are able to learn complex, non-linear patterns but usually need more and less interpretable data. Comparative experiments show that the simpler models often perform competitively with respect to more complex alternatives especially using small datasets. To the current use case, i.e. predicting interest in particular software features with the help of public data, an ensemble of classical methods can offer the right level of accuracy and interpretability. Current research on scalable forecasting has shown that hybrid systems that incorporate statistical rigour with practical implementation strategies are showing excellent results.

\subsection{Sentiment Analysis and User Feedback}
The feedback provided by users is one of the most useful sources of information about user sentiment, as it is taken in the form of reviews in the app shops and discussions in social media. NLP methods allow studying extensive review collections in order to identify feature mentions, pain points, and satisfaction patterns. Sentiment analysis gives extra information about what is discussed by the users in addition to their emotional inclination towards certain features. Other literature in this field has concentrated more on the existing user perceptions about existing features. The study is unique in that it tries to forecast future demand and not just quantify the current sentiment. Nevertheless, temporal patterns in sentiment have the potential to mean demand trends, which is why sentiment analysis should be implemented into a forecasting model.

\subsection{Research Gap and Contribution}
The research void that the study will fill is that there are no unified frameworks relating to the design of SaaS products in terms of forecasting, prioritisation, and sentiment analysis. Traditionally forecasting research has been concerned with the physical products. The systems of priorities are still qualitative. The work of sentiment analysis is mainly retrospective instead of being predictive. This paper forms a framework that incorporates these lines of research. Using a combination of forecasting methods, testing on empirical trends in features, and only using publicly available data, this study shows a practical method that can be used by a product team that does not have a large analytical infrastructure. It is methodological but not algorithmic in nature of contribution---larger teams are made available to small teams of data-driven decision-making capabilities that would otherwise require large analytical resources.
